\chapter{Kernenergie: splijting, fusie, $E=mc^2$}\label{ch:g12:nuclear-energy}

Eén brandstoftablet ter grootte van een gummetje dekt de elektriciteit van een
huishouden voor een jaar; dezelfde dienst met steenkool kost drie ton. Boven
ons hoofd betaalt de Zon het daglicht door te verdwijnen --- vier miljoen ton
van zichzelf per seconde. Dit hoofdstuk opent het rekeningenboek achter beide
koopjes: massa is zelf \omterm{def:g10:energy-conservation:energy}{energie}, en in de kernen
(\cref{ch:g11:nucleus-radioactivity}) wordt ze verzilverd.

\section{Massa is energie}

\begin{theorem}[Equivalentie van massa en energie]\label{thm:g12:nuclear-energy:emc2}\index{massa-energie-equivalentie}
Een lichaam met massa $m$ bevat, alleen al door te bestaan, de \omterm{def:g10:energy-conservation:energy}{energie}
\[
E = m c^2, \qquad c = \qty{3.00e8}{m/s}.
\]
Telkens wanneer een systeem massa $\Delta m$ verliest, geeft het de \omterm{def:g10:energy-conservation:energy}{energie}
$\Delta m\, c^2$ vrij --- wat het mechanisme ook is.
\end{theorem}
\admitted

\begin{remark}[Einstein, 1905]\label{rem:g12:nuclear-energy:einstein}
Dit is de betrekking van Einstein, hier toegegeven; waar ze vandaan komt is de
zaak van het hoofdstuk over de speciale relativiteit dat dit jaar afsluit
(\cref{ch:g12:special-relativity}). Let op de koers:
$c^2 = \qty{9.00e16}{J/kg}$ --- een afrondingsfout aan massa is een fortuin
aan \omterm{def:g10:energy-conservation:energy}{energie}.
\end{remark}

\begin{definition}[Elektronvolt en atomaire massa-eenheid]\label{def:g12:nuclear-energy:units}
Twee eenheden op maat van de kern. De
\emph{elektronvolt}\index{elektronvolt}, de \omterm{def:g10:energy-conservation:energy}{energie} die een elektron wint door
\qty{1}{V} te doorlopen (\cref{ch:g11:circuits-and-power}):
$\qty{1}{eV} = \qty{1.60e-19}{J}$, $\qty{1}{MeV} = \qty{1.60e-13}{J}$ ---
chemische bindingen verhandelen enkele \unit{eV}, kernreacties
mega-elektronvolt. En de
\emph{atomaire massa-eenheid}\index{atomaire massa-eenheid}, een twaalfde van
de massa van een koolstof-12-atoom:
$\qty{1}{u} = \qty{1.66054e-27}{kg}$, met energie-equivalent
\qty{931.5}{MeV} --- $\qty{1}{u} = \qty{931.5}{MeV}/c^2$.
\end{definition}

\begin{example}[Het slapende fortuin]\label{ex:g12:nuclear-energy:fortune}
Eén \omterm{def:g10:orders-of-magnitude:unit}{kilogram} van om het even wat is, volledig omgezet, $\qty{9.00e16}{J}$ ---
drie jaar productie van een grote centrale (\qty{1}{GW}). Niemand merkte het
vóór 1905, omdat de gewone natuurkunde massa nauwelijks aanraakt: een \omterm{def:g10:orders-of-magnitude:unit}{kilogram}
steenkool verbranden (\qty{3.0e7}{J}) maakt hem een derde microgram lichter.
Kernen, zullen we zien, verschuiven één deel op duizend: nog altijd klein, maar
een miljoen keer de chemie.
\end{example}

\section{Het massadefect}

\begin{definition}[Massadefect en bindingsenergie]\label{def:g12:nuclear-energy:defect}
Weeg een kern ${}^{A}_{Z}\mathrm{X}$ (massa $m$), en weeg daarna haar $Z$
protonen en $A - Z$ neutronen apart: de delen zijn \emph{zwaarder} dan het
geheel. Het verschil
\[
\Delta m = Z\,m_p + (A - Z)\,m_n - m > 0,
\qquad m_p = \qty{1.00728}{u},\; m_n = \qty{1.00866}{u},
\]
is het \emph{massadefect}\index{massadefect}; de
\emph{bindingsenergie}\index{bindingsenergie} $E_b = \Delta m\,c^2$ is de
\omterm{def:g10:energy-conservation:energy}{energie} die nodig is om de kern uit elkaar te trekken tot vrije \omterm{def:g11:fundamental-interactions:ladder}{nucleonen} ---
en evengoed de \omterm{def:g10:energy-conservation:energy}{energie} die vrijkomt wanneer ze wordt samengesteld. Een gebonden
kern ligt \emph{onder} haar delen.
\end{definition}

\begin{omfigure}
\begin{tikzpicture}[font=\small]
  \draw[thick] (0,2.1) -- (2.7,2.1) node[right] {$Z$ protonen $+$ $(A-Z)$ neutronen, vrij};
  \foreach \x in {0.3,0.9,1.5,2.1} \fill[omThm!80] (\x,2.4) circle (2.4pt);
  \foreach \x in {0.6,1.2,1.8,2.4} \fill[omExo] (\x,2.4) circle (2.4pt);
  \draw[thick, omDef] (0,0) -- (2.7,0) node[right] {kern, massa $m$};
  \draw[-Stealth, thick] (1.35,2.1) -- (1.35,0) node[midway, right] {$E_b = \Delta m\,c^2$};
\end{tikzpicture}

{\small Een energieladder: de gebonden kern ligt $E_b$ onder haar gescheiden
\omterm{def:g11:fundamental-interactions:ladder}{nucleonen} --- haar samenstellen geeft $E_b$ vrij, haar ontmantelen kost het.}
\end{omfigure}

\begin{example}[Helium-4]\label{ex:g12:nuclear-energy:helium}
De kern ${}^{4}_{2}\mathrm{He}$ heeft $m = \qty{4.00151}{u}$; haar delen:
$2 \times 1.00728 + 2 \times 1.00866 = \qty{4.03188}{u}$. Dus
$\Delta m = \qty{0.03037}{u}$ en $E_b = 0.03037 \times 931.5 \approx
\qty{28.3}{MeV}$: de kern weegt $0.75\%$ minder dan haar delen --- het ``één
deel op duizend'' van \cref{ex:g12:nuclear-energy:fortune}.
\end{example}

\begin{method}[Energiebalans van een kernreactie]\label{met:g12:nuclear-energy:balance}
\begin{enumerate}
  \item Laat de vergelijking kloppen: $A$ en $Z$ blijven in totaal
        behouden (\cref{ch:g11:nucleus-radioactivity}).
  \item Tel de massa's ervóór op, daarna die erna, in \unit{u} --- houd alle
        vijf de decimalen aan.
  \item $\Delta m = m_{\text{ervoor}} - m_{\text{erna}}$, en daarna
        $E = \Delta m\,(\text{in u}) \times \qty{931.5}{MeV}$, die als
        \omterm{def:g11:mechanical-energy:kinetic}{kinetische energie} van de producten wordt meegevoerd; negatief betekent
        dat de reactie betaald moet worden.
\end{enumerate}
\end{method}

\section{De kromme die het heelal bestuurt}

\begin{definition}[Bindingsenergie per nucleon]\label{def:g12:nuclear-energy:pernucleon}
Samenhang vergelijk je eerlijk met de \emph{bindingsenergie per
nucleon}\index{bindingsenergie per nucleon}, $E_b/A$, de gemiddelde prijs om
één \omterm{def:g11:fundamental-interactions:ladder}{nucleon} los te maken: $28.3/4 \approx \qty{7.1}{MeV}$ voor helium-4.
\end{definition}

\begin{proposition}[De ijzerpiek, en twee wegen omlaag]\label{prop:g12:nuclear-energy:ironpeak}\index{ijzerpiek}
Uitgezet tegen $A$ klimt de \omterm{def:g12:nuclear-energy:pernucleon}{bindingsenergie per nucleon} steil door de lichte
kernen, piekt ze bij ijzer, ${}^{56}_{26}\mathrm{Fe}$, op ongeveer
\qty{8.8}{MeV} per \omterm{def:g11:fundamental-interactions:ladder}{nucleon}, en glijdt ze daarna zacht omlaag naar zo'n
\qty{7.6}{MeV} bij uranium: ijzer is de sterkst gebonden kern in de natuur.
Een reactie geeft precies dan \omterm{def:g10:energy-conservation:energy}{energie} vrij wanneer haar producten \emph{hoger}
op de kromme liggen; twee tegengestelde strategieën lonen dus allebei, en
allebei lopen ze naar ijzer toe --- de zwaarste kernen \emph{splijten}, en de
lichtste \emph{samensmelten}.
\end{proposition}
\admitted

\begin{remark}[Waarom een piek]\label{rem:g12:nuclear-energy:whypeak}
Het touwtrekken van \cref{ch:g11:fundamental-interactions}: de sterke lijm
bindt alleen buren die elkaar raken, terwijl de coulombafstoting de hele kern
overspant --- kleine kernen zijn vooral oppervlak, grote vooral afstoting, en
ijzer is het compromis. De eerlijke boekhouding woont in de universitaire
volumes.
\end{remark}

\begin{omfigure}
\begin{tikzpicture}
\begin{axis}[omaxis, width=9.2cm, height=5.6cm,
    xlabel={massagetal $A$}, ylabel={$E_b/A$ (\unit{MeV})},
    xmin=0, xmax=248, ymin=0, ymax=10.4, ytick={0,2,4,6,8}, xtick={4,56,120,180,235}]
  \addplot[omDef, thick, smooth] coordinates {(1,0) (2,1.1) (3,2.8) (6,5.3) (9,6.5)
    (12,7.7) (20,8.0) (40,8.6) (56,8.8) (90,8.7) (120,8.5) (150,8.3) (190,7.9) (235,7.6)};
  \addplot[only marks, mark=*, mark size=1.6pt, omThm] coordinates {(4,7.1) (56,8.8) (235,7.6)};
  \node[font=\small, right] at (axis cs:8,7.1) {$^{4}$He};
  \node[font=\small, above] at (axis cs:56,8.9) {$^{56}$Fe};
  \node[font=\small, above] at (axis cs:235,7.75) {$^{235}$U};
  \draw[-Stealth, omProp, thick] (axis cs:16,3.4) -- (axis cs:44,7.6)
    node[midway, below right, font=\small] {kernfusie};
  \draw[-Stealth, omProp, thick] (axis cs:222,6.9) -- (axis cs:130,7.8)
    node[midway, below, font=\small] {kernsplijting};
\end{axis}
\end{tikzpicture}

{\small \omterm{def:g12:nuclear-energy:pernucleon}{Bindingsenergie per nucleon}: steile klim, \omterm{prop:g12:nuclear-energy:ironpeak}{ijzerpiek}, lange glijbaan
--- \omterm{def:g12:nuclear-energy:fusion}{kernfusie} klimt van links, \omterm{def:g12:nuclear-energy:fission}{kernsplijting} van rechts, en allebei geven ze
\omterm{def:g10:energy-conservation:energy}{energie} vrij.}
\end{omfigure}

\section{Splijting en de kettingreactie}

\begin{definition}[Kernsplijting]\label{def:g12:nuclear-energy:fission}
\emph{Kernsplijting}\index{kernsplijting} is het uiteenvallen van een zware
kern in twee middelgrote brokstukken. Uranium-235 is
\emph{splijtbaar}\index{splijtbaar}: een traag neutron dat wordt opgenomen,
laat de kern zo opgewonden achter dat ze in tweeën scheurt en daarbij twee of
drie verse neutronen uitspuwt ---
${}^{1}_{0}\mathrm{n} + {}^{235}_{92}\mathrm{U} \to \text{twee brokstukken} +
2\text{--}3\;{}^{1}_{0}\mathrm{n}$, met een paar dat van splijting tot
splijting wisselt.
\end{definition}

\begin{example}[Eén splijting, gewogen]\label{ex:g12:nuclear-energy:fissionbalance}
Eén veelvoorkomende route, volgens \cref{met:g12:nuclear-energy:balance}:
\[
{}^{1}_{0}\mathrm{n} + {}^{235}_{92}\mathrm{U} \longrightarrow
{}^{94}_{38}\mathrm{Sr} + {}^{140}_{54}\mathrm{Xe} + 2\,{}^{1}_{0}\mathrm{n}.
\]
Ervóór: $1.00866 + 234.99350 = \qty{236.00216}{u}$. Erna:
$93.89450 + 139.89200 + 2 \times 1.00866 = \qty{235.80382}{u}$. Dus is
$\Delta m = \qty{0.19834}{u}$ en $E \approx \qty{185}{MeV}$; tel je de latere
vervallen van de radioactieve brokstukken mee
(\cref{ch:g12:radioactive-decay}), dan is elke \omterm{def:g12:nuclear-energy:fission}{kernsplijting} ongeveer
\qty{200}{MeV} waard --- terwijl één \omterm{def:g11:fundamental-interactions:ladder}{atoom} steenkool voor enkele \unit{eV} te
koop is.
\end{example}

\begin{definition}[Kettingreactie en kritische massa]\label{def:g12:nuclear-energy:chain}
Elke \omterm{def:g12:nuclear-energy:fission}{kernsplijting} wordt door één neutron ontstoken en bevrijdt er twee of
drie: in een klomp die groot genoeg is, treffen die andere kernen en voedt de
reactie zichzelf --- een \emph{kettingreactie}\index{kettingreactie}. In een
kleine klomp ontsnappen de meeste neutronen eerst door het oppervlak; de
kleinste hoeveelheid die de ketting in stand houdt, is de \emph{kritische
massa}\index{kritische massa} --- daaronder dooft de ketting, daarboven
vermenigvuldigt elke generatie zich: het beginsel van de bom.
\end{definition}

\begin{omfigure}
\begin{tikzpicture}[font=\small, scale=0.9]
  \draw[-Stealth, thick] (-1.4,0) -- (-0.4,0) node[pos=0.15, above] {n};
  \fill[omDef!55] (0,0) circle (11pt); \node at (0,0) {$^{235}$U};
  \fill[omExo] (0.75,0.85) circle (4.5pt) node[above=3pt] {brokstukken};
  \fill[omExo] (0.75,-0.85) circle (4.5pt);
  \foreach \y in {1.8,0,-1.8} {\draw[-Stealth] (0.5,0) -- (2.5,\y);
    \fill[omDef!55] (3,\y) circle (11pt); \node at (3,\y) {$^{235}$U};
    \foreach \a in {28,0,-28} \draw[-Stealth] (3.5,\y) -- ++(\a:1.1);}
\end{tikzpicture}

{\small Een \omterm{def:g12:nuclear-energy:chain}{kettingreactie}: één neutron splijt één kern, en de drie bevrijde
neutronen splijten er drie --- een reactor laat er precies één doorgaan.}
\end{omfigure}

\begin{definition}[De getemde reactor]\label{def:g12:nuclear-energy:reactor}
Een energiereactor houdt de ketting op precies één neutron per \omterm{def:g12:nuclear-energy:fission}{kernsplijting}.
De brandstof: uraniumtabletten die bescheiden in uranium-235 zijn verrijkt. De
\emph{moderator}\index{moderator} (meestal gewoon water) vertraagt de snelle
splijtingsneutronen door botsingen --- trage neutronen zijn veel beter in het
ontsteken van de volgende splijting; de
\emph{regelstaven}\index{regelstaaf} (boor, cadmium) verslinden neutronen ---
ingeschoven dooft de ketting, uitgetrokken versnelt ze. De opwinding van de
brokstukken wordt warmte, daarna stoom, daarna elektriciteit
(\cref{ch:g10:energy-conservation}).
\end{definition}

\begin{example}[Een peperkorrel tegen een vrachtwagen]\label{ex:g12:nuclear-energy:gram}
Eén gram uranium-235 bevat
$N = \num{1.0e-3} / (235 \times \num{1.66054e-27}) \approx \num{2.56e21}$
kernen; bij $\qty{200}{MeV} = \qty{3.2e-11}{J}$ elk levert ze allemaal
splijten ongeveer \qty{8.2e10}{J} op --- de \omterm{def:g10:energy-conservation:forms}{chemische energie} van
$\num{8.2e10} / \num{3.0e7} \approx \qty{2.7e3}{kg}$ steenkool: een
peperkorrel tegen drie ton.
\end{example}

\section{Fusie: het eigen vuur van de Zon}

\begin{definition}[Kernfusie]\label{def:g12:nuclear-energy:fusion}
\emph{Kernfusie}\index{kernfusie} is het samensmelten van twee lichte kernen
tot een zwaardere, steviger gebonden kern. De best bereikbare reactie huwt
deuterium en tritium, de zware \omterm{def:g11:nucleus-radioactivity:notation}{isotopen} van waterstof
(\cref{ch:g11:nucleus-radioactivity}):
\[
{}^{2}_{1}\mathrm{H} + {}^{3}_{1}\mathrm{H} \longrightarrow
{}^{4}_{2}\mathrm{He} + {}^{1}_{0}\mathrm{n}.
\]
\end{definition}

\begin{example}[D--T, gewogen]\label{ex:g12:nuclear-energy:dt}
Ervóór: $2.01355 + 3.01550 = \qty{5.02905}{u}$; erna:
$4.00151 + 1.00866 = \qty{5.01017}{u}$; dus $\Delta m = \qty{0.01888}{u}$ en
$E \approx \qty{17.6}{MeV}$, waarvan vier vijfde op het neutron --- minder dan
een \omterm{def:g12:nuclear-energy:fission}{kernsplijting}, maar uit vijf \omterm{def:g11:fundamental-interactions:ladder}{nucleonen} in plaats van $236$: ongeveer
\qty{3.4e14}{J} per \omterm{def:g10:orders-of-magnitude:unit}{kilogram} brandstof, vier keer de splijting. Helium-4
(\cref{ex:g12:nuclear-energy:helium}) is opnieuw de as.
\end{example}

\begin{definition}[De coulombbarrière]\label{def:g12:nuclear-energy:barrier}
Twee kernen zijn allebei positief: om elkaar te raken moeten ze eerst de
heuvel van hun elektrische afstoting op
(\cref{ch:g11:electric-gravitational-fields}) --- de
\emph{coulombbarrière}\index{coulombbarrière}. Pas rond \qty{1e8}{K} dragen de
botsingen de \omterm{def:g10:energy-conservation:energy}{energie} om eroverheen te komen; de brandstof is dan een
\emph{plasma}\index{plasma} van kale kernen en elektronen. \omterm{def:g12:nuclear-energy:fission}{Kernsplijting}
betaalt zo'n tol niet: haar ontsteker, het neutron, is neutraal en werkt koud.
\end{definition}

\begin{example}[De Zon draait erop]\label{ex:g12:nuclear-energy:sun}
De kern van de Zon, op $\qty{1.5e7}{K}$, smelt gewone waterstof stap voor stap
samen --- de \emph{proton-protonketen}, met als netto-effect
$4\,{}^{1}_{1}\mathrm{H} \to {}^{4}_{2}\mathrm{He} + 2\,{}^{0}_{+1}\mathrm{e}
+ \text{straling}$, ongeveer \qty{26}{MeV} per helium: $0.7\%$ van de massa
van de waterstof vertrekt als \omterm{def:g10:energy-conservation:energy}{energie}. Het zonlicht waarvan we het \omterm{def:g12:sound-acoustics:timbre}{spectrum} in
\cref{ch:g10:light-spectra} lezen, is die massa, acht minuten te laat
aangekomen; elke ster schijnt door de kromme naar ijzer af te lopen.
\end{example}

\begin{remark}[ITER: een ster op fles]\label{rem:g12:nuclear-energy:iter}
Geen enkele vaste stof houdt een \omterm{def:g12:nuclear-energy:barrier}{plasma} van \qty{1.5e8}{K} vast; een
\emph{tokamak} kooit het in een \omterm{def:g11:magnetic-fields:bfield}{magnetisch veld}
(\cref{ch:g11:magnetic-fields}), een donut van \omterm{def:g11:electric-gravitational-fields:lines}{veldlijnen} waarlangs de geladen
deeltjes spiralen zonder de wand te raken; ITER, gebouwd om tien keer zijn
verwarmingsvermogen vrij te geven, is de huidige stap. Breek de opsluiting en
het \omterm{def:g12:nuclear-energy:barrier}{plasma} koelt af en stopt binnen seconden: geen \omterm{def:g12:nuclear-energy:chain}{kettingreactie}, geen
\omterm{def:g12:nuclear-energy:chain}{kritische massa} --- het moeilijke is niet het vuur doven maar het aan houden.
\end{remark}

\begin{example}[De energieladder, per kilogram]\label{ex:g12:nuclear-energy:ladder}
\omterm{def:g10:energy-conservation:energy}{Energie} per \omterm{def:g10:orders-of-magnitude:unit}{kilogram} brandstof: steenkool, \qty{3.0e7}{J}; uranium-235 door
\omterm{def:g12:nuclear-energy:fission}{kernsplijting}, \qty{8.2e13}{J}; deuterium--tritium door \omterm{def:g12:nuclear-energy:fusion}{kernfusie},
\qty{3.4e14}{J}; volledige omzetting ($E = mc^2$), \qty{9.0e16}{J}. Zes
\omterm{def:g10:orders-of-magnitude:oom}{ordegrootten} scheiden de chemie van de kern --- de reden dat een reactor per
vrachtwagen wordt bevoorraad en een kolencentrale per goederentrein.
\end{example}

\section{Oefeningen}

\begin{exercise}[$\star$]\label{exo:g12:nuclear-energy:1}
Reken om: (a) \qty{1}{eV} en \qty{1}{MeV} naar \omterm{def:g11:work-of-force:work}{joule}; (b) de \qty{200}{MeV}
van één \omterm{def:g12:nuclear-energy:fission}{kernsplijting} naar \omterm{def:g11:work-of-force:work}{joule}; (c) ga uit
$\qty{1}{u} = \qty{1.66054e-27}{kg}$ na dat haar energie-equivalent dicht bij
\qty{931.5}{MeV} ligt.
\end{exercise}

\begin{exercise}[$\star$]\label{exo:g12:nuclear-energy:2}
Een suikerklontje heeft een massa van \qty{6.0}{g}: bereken zijn volle waarde
$E = mc^2$. Een huishouden verbruikt $\qty{10}{kW.h} = \qty{3.6e7}{J}$ per
dag --- hoeveel jaar zou het klontje, volledig omgezet, het van stroom kunnen
voorzien?
\end{exercise}

\begin{exercise}[$\star$]\label{exo:g12:nuclear-energy:3}
De koolstof-12-kern heeft massa \qty{11.99671}{u}: bereken haar \omterm{def:g12:nuclear-energy:defect}{massadefect},
haar \omterm{def:g12:nuclear-energy:defect}{bindingsenergie} in \unit{MeV} en haar \omterm{def:g12:nuclear-energy:pernucleon}{bindingsenergie per nucleon}
($m_p = \qty{1.00728}{u}$, $m_n = \qty{1.00866}{u}$).
\end{exercise}

\begin{exercise}[$\star$]\label{exo:g12:nuclear-energy:4}
Vul elke \omterm{def:g12:nuclear-energy:fission}{kernsplijting} aan met $A$, $Z$ en het element ($Z = 36$: krypton,
$Z = 52$: telluur):
(a) ${}^{1}_{0}\mathrm{n} + {}^{235}_{92}\mathrm{U} \to
{}^{144}_{56}\mathrm{Ba} + {}^{A}_{Z}\mathrm{X} + 3\,{}^{1}_{0}\mathrm{n}$;
(b) ${}^{1}_{0}\mathrm{n} + {}^{235}_{92}\mathrm{U} \to
{}^{97}_{40}\mathrm{Zr} + {}^{A}_{Z}\mathrm{Y} + 2\,{}^{1}_{0}\mathrm{n}$.
\end{exercise}

\begin{exercise}[$\star$]\label{exo:g12:nuclear-energy:5}
Lees de kromme van \cref{prop:g12:nuclear-energy:ironpeak}: schat $E_b/A$ voor
helium-4, koolstof-12, ijzer-56 en uranium-235. Welke is het sterkst gebonden?
Welke zou \omterm{def:g10:energy-conservation:energy}{energie} kunnen vrijgeven door \omterm{def:g12:nuclear-energy:fission}{kernsplijting}, en welke door
\omterm{def:g12:nuclear-energy:fusion}{kernfusie}?
\end{exercise}

\begin{exercise}[$\star\star$]\label{exo:g12:nuclear-energy:6}
Het deuteron ${}^{2}_{1}\mathrm{H}$ heeft massa \qty{2.01355}{u}: bereken zijn
\omterm{def:g12:nuclear-energy:defect}{bindingsenergie} in \unit{MeV}. Een chemische binding houdt enkele \unit{eV}
vast --- met welke factor verslaat de ``kernbinding'' haar?
\end{exercise}

\begin{exercise}[$\star\star$]\label{exo:g12:nuclear-energy:7}
Weeg de \omterm{def:g12:nuclear-energy:fission}{kernsplijting} ${}^{1}_{0}\mathrm{n} + {}^{235}_{92}\mathrm{U} \to
{}^{144}_{56}\mathrm{Ba} + {}^{89}_{36}\mathrm{Kr} + 3\,{}^{1}_{0}\mathrm{n}$:
massa's \qty{234.99350}{u} (U), \qty{143.89220}{u} (Ba), \qty{88.89790}{u}
(Kr), \qty{1.00866}{u} (n). Hoeveel \omterm{def:g10:energy-conservation:energy}{energie} komt er vrij, in \unit{MeV} en in
\omterm{def:g11:work-of-force:work}{joule}?
\end{exercise}

\begin{exercise}[$\star\star$]\label{exo:g12:nuclear-energy:8}
Nog een \omterm{def:g12:nuclear-energy:fusion}{kernfusie}: ${}^{2}_{1}\mathrm{H} + {}^{2}_{1}\mathrm{H} \to
{}^{3}_{2}\mathrm{He} + {}^{1}_{0}\mathrm{n}$, met
$m({}^{3}_{2}\mathrm{He}) = \qty{3.01493}{u}$. Bereken de vrijgekomen \omterm{def:g10:energy-conservation:energy}{energie}.
Waarom zoveel minder dan D--T? (Kijk waar elk product op de kromme ligt.)
\end{exercise}

\begin{exercise}[$\star\star$]\label{exo:g12:nuclear-energy:9}
Leg uit, telkens in een paar regels: (a) waarom \omterm{def:g12:nuclear-energy:fission}{kernsplijting} geen verwarming
nodig heeft terwijl \omterm{def:g12:nuclear-energy:fusion}{kernfusie} zo'n \qty{1e8}{K} vraagt; (b) hoe de Zon het met
``maar'' \qty{1.5e7}{K} redt (denk aan de omvang van de menigte en aan haar
snelste leden); (c) waarom een fusiecentrale niet zoals een \omterm{def:g12:nuclear-energy:chain}{kettingreactie} op
hol kan slaan.
\end{exercise}

\begin{exercise}[$\star\star$]\label{exo:g12:nuclear-energy:10}
De nettoreactie van de Zon maakt van vier protonen ($m_p = \qty{1.00728}{u}$)
één helium-4-kern (\qty{4.00151}{u}): bereken de \omterm{def:g10:energy-conservation:energy}{energie} per gemaakt helium,
in \unit{MeV}, en de fractie van de beginmassa die wordt omgezet.
\end{exercise}

\begin{exercise}[$\star\star$]\label{exo:g12:nuclear-energy:11}
In een reactor met water als \omterm{def:g12:nuclear-energy:reactor}{moderator}: (a) wat doet de \omterm{def:g12:nuclear-energy:reactor}{moderator}, en waarom
heeft de ketting hem nodig? (b) en de \omterm{def:g12:nuclear-energy:reactor}{regelstaven}? (c) het koelwater \emph{is}
de \omterm{def:g12:nuclear-energy:reactor}{moderator}: waarom neigt het verlies ervan de ketting eerder te smoren dan
te versnellen?
\end{exercise}

\begin{exercise}[$\star\star\star$]\label{exo:g12:nuclear-energy:12}
Een ongecontroleerde ketting verdubbelt bij elke generatie. Hoeveel generaties
duurt het vanaf één \omterm{def:g12:nuclear-energy:fission}{kernsplijting} ongeveer tot de $\num{2.56e21}$ kernen van
één gram uranium-235 zijn gespleten ($2^{10} \approx 10^{3}$)? Hoe lang is dat
bij ruwweg \qty{10}{ns} per generatie --- en waarom wordt kriticiteit dan zo
zorgvuldig bewaakt?
\end{exercise}

\begin{exercise}[$\star\star\star$]\label{exo:g12:nuclear-energy:13}
Zouden we \omterm{def:g10:energy-conservation:energy}{energie} kunnen winnen door ijzer te splijten? Weeg
${}^{56}_{26}\mathrm{Fe} \to 2\,{}^{28}_{14}\mathrm{Si}$, met
$m(\mathrm{Fe}) = \qty{55.92066}{u}$ en $m(\mathrm{Si}) = \qty{27.96924}{u}$.
Bereken $\Delta m$ en concludeer, met de kromme in de hand, waarom ijzer
kernas is: brandstof voor niets.
\end{exercise}

\begin{exercise}[$\star\star\star$]\label{exo:g12:nuclear-energy:14}
Bereken uit $m({}^{56}_{26}\mathrm{Fe}) = \qty{55.92066}{u}$ en
$m({}^{235}_{92}\mathrm{U}) = \qty{234.99350}{u}$ de $E_b/A$ van beide
($m_p = \qty{1.00728}{u}$, $m_n = \qty{1.00866}{u}$). Splijtingsbrokstukken
liggen rond \qty{8.5}{MeV} per \omterm{def:g11:fundamental-interactions:ladder}{nucleon}: \emph{schat} uit de twee niveaus de
\omterm{def:g10:energy-conservation:energy}{energie} van één \omterm{def:g12:nuclear-energy:fission}{kernsplijting}, en vergelijk met
\cref{ex:g12:nuclear-energy:fissionbalance}.
\end{exercise}

\begin{exercise}[$\star\star\star$]\label{exo:g12:nuclear-energy:15}
ITER mikt op \qty{500}{MW} fusievermogen. Bepaal bij \qty{17.6}{MeV} per
D--T-reactie (\qty{5.029}{u} verbruikte brandstof) het aantal reacties per
seconde, en daarna de verbrande brandstof per dag. Een kolenoven van
\qty{500}{MW} verstookt ongeveer \qty{1.4e6}{kg} per dag: vergelijk.
\end{exercise}

\section{Probleem: Een stad van stroom voorzien}

\begin{problem}[{Weekendopgave --- een stad van stroom voorzien: één gestage
gigawatt voor een half miljoen mensen, op vier manieren gekocht ---
uraniumtabletten, treinladingen steenkool, een zwembad zeewater, de eigen
substantie van de Zon --- met onder alle vier dezelfde verdwenen kilogram
massa}]
\label{pb:g12:nuclear-energy:1}
Een stad van $500\,000$ inwoners trekt een gestaag \omterm{prop:g11:circuits-and-power:power}{elektrisch vermogen}
$P = \qty{1.0}{GW}$; haar centrales zetten warmte, wat de brandstof ook is, met
een \omterm{def:g10:energy-conservation:efficiency}{rendement} van $33\%$ in elektriciteit om
(\cref{ch:g10:energy-conservation}). Gegevens: één \omterm{def:g12:nuclear-energy:fission}{kernsplijting} van
uranium-235, \qty{200}{MeV} in totaal; steenkool, \qty{3.0e7}{J/kg}; één
D--T-fusie, \qty{17.6}{MeV} uit \qty{5.029}{u} brandstof;
$\qty{1}{u} = \qty{1.66054e-27}{kg}$; $\qty{1}{eV} = \qty{1.60e-19}{J}$; één
jaar $= \qty{3.16e7}{s}$; de Zon: \qty{3.9e26}{W}, \qty{2.0e30}{kg}; een
kubieke \omterm{def:g10:orders-of-magnitude:unit}{meter} zeewater bevat ongeveer \qty{33}{g} deuterium.

\medskip
\noindent\textbf{Deel I --- De splijtingscentrale.}
\begin{enumerate}
  \item Welk thermisch \omterm{def:g11:work-of-force:power}{vermogen} moet de reactor leveren?
  \item \omterm{def:g10:pressure:pressure}{Druk} de \omterm{def:g10:energy-conservation:energy}{energie} van één \omterm{def:g12:nuclear-energy:fission}{kernsplijting} in \omterm{def:g11:work-of-force:work}{joule} uit.
  \item Hoeveel \omterm{def:g12:nuclear-energy:fission}{kernsplijtingen} per seconde houden de stad draaiende?
  \item Wat is de massa van één uranium-235-kern ($A = 235$), en hoeveel
        uranium-235 wordt er per seconde verbruikt?
  \item Leid daaruit het verbruik van uranium-235 in één jaar af, in \omterm{def:g10:orders-of-magnitude:unit}{kilogram}.
  \item De brandstof is tot $4.0\%$ uranium-235 verrijkt: welke massa vraagt
        een jaar, en past die op één vrachtwagen?
\end{enumerate}

\medskip
\noindent\textbf{Deel II --- Het kolengrootboek.}
\begin{enumerate}[resume]
  \item Bereken de \omterm{def:g10:energy-conservation:forms}{thermische energie} die de stad in één jaar verbruikt.
  \item Welke massa steenkool levert die?
  \item Bereken de verhouding (massa steenkool)/(massa uranium-235) voor één
        jaar, en toets die aan de ladder van
        \cref{ex:g12:nuclear-energy:ladder}.
  \item Een goederentrein sleept \qty{3.0e6}{kg}: hoeveel treinen per jaar, en
        per dag? Beschrijf in één zin de twee bevoorradingslijnen.
\end{enumerate}

\medskip
\noindent\textbf{Deel III --- De fusiedroom.}
\begin{enumerate}[resume]
  \item \omterm{def:g10:pressure:pressure}{Druk} de \omterm{def:g10:energy-conservation:energy}{energie} van één D--T-fusie in \omterm{def:g11:work-of-force:work}{joule} uit.
  \item Bereken de vrijgekomen \omterm{def:g10:energy-conservation:energy}{energie} per \omterm{def:g10:orders-of-magnitude:unit}{kilogram} D--T-brandstof; vergelijk
        met de ladder.
  \item Welke massa D--T-brandstof dekt het jaar van de stad? En welke massa
        daarvan is deuterium ($2.014$ van de \qty{5.029}{u})?
  \item Welk volume zeewater bevat dat deuterium? Vergelijk met een olympisch
        zwembad, ongeveer \qty{2.5e3}{m^3}. (Tritium wordt uit lithium
        gekweekt.)
  \item Waarom draagt deze centrale geen \omterm{def:g12:nuclear-energy:chain}{kritische massa} en geen
        \omterm{def:g12:nuclear-energy:chain}{kettingreactie} om bang voor te zijn? Wat is dan het moeilijke?
\end{enumerate}

\medskip
\noindent\textbf{Deel IV --- De ster die in massa betaalt.}
\begin{enumerate}[resume]
  \item Welke massa zet de Zon volgens $E = mc^2$ per seconde om?
  \item Welke massa per seconde, en per jaar, moet \emph{elke} thermische bron
        van \qty{3.0}{GW} omzetten? Ga na dat alle drie de brandstoffen
        diezelfde massa afstaan.
  \item Welke totale massa heeft de Zon in haar $\num{4.5e9}$ jaar
        weggestraald, en welke fractie van de Zon is dat?
  \item \omterm{def:g12:nuclear-energy:fusion}{Kernfusie} kan ongeveer $0.07\%$ van de massa van de Zon aanspreken
        (het brandbare deel van haar kern): schat haar totale levensduur als
        ster, en wat daarvan overblijft.
  \item Sluit de audit in drie zinnen af: één jaar van de stad in tonnen
        steenkool, \omterm{def:g10:orders-of-magnitude:unit}{kilogram} uranium-235 en \omterm{def:g10:orders-of-magnitude:unit}{kilogram} D--T-brandstof --- en de
        \omterm{def:g10:orders-of-magnitude:unit}{kilogram} massa die in elk scenario het licht van de stad is geworden.
\end{enumerate}
\end{problem}
